% ============================================
% PLANTILLA PARA TESIS UNTELS
% ============================================
% Uso: Tesis de pregrado - Ingeniería Electrónica
% Estilo: Formato institucional UNTELS
% ============================================

\documentclass[12pt, letterpaper, oneside]{book}

% ============================================
% PAQUETES REQUERIDOS
% ============================================
\usepackage[utf8]{inputenc}
\usepackage[spanish]{babel}
\usepackage{csquotes}      % requerido por biblatex con babel
\usepackage{amsmath, amssymb, amsfonts}
\usepackage{graphicx}
\usepackage{float}
\usepackage{booktabs}
\usepackage{longtable}
\usepackage{hyperref}
\usepackage{geometry}
\usepackage{fancyhdr}
\usepackage{titlesec}
\usepackage{setspace}
\usepackage{appendix}
\usepackage{listings}
% Permite tildes y ñ dentro de lstlisting con inputenc utf8
\lstset{
    inputencoding=utf8, extendedchars=true,
    literate={á}{{\'a}}1 {é}{{\'e}}1 {í}{{\'i}}1 {ó}{{\'o}}1 {ú}{{\'u}}1
             {Á}{{\'A}}1 {É}{{\'E}}1 {Í}{{\'I}}1 {Ó}{{\'O}}1 {Ú}{{\'U}}1
             {ñ}{{\~n}}1 {Ñ}{{\~N}}1
}
\usepackage{xcolor}
\usepackage{multirow}
\usepackage{array}
\usepackage{rotating}

% ============================================
% BIBLIOGRAFÍA: APA 7.ª EDICIÓN (biblatex-apa + biber)
% ============================================
% Compilar: pdflatex -> biber -> pdflatex -> pdflatex  (o: latexmk -pdf tesis_untels.tex)
% ¡Use BIBER, no BibTeX! En Overleaf: Menu > Compiler pdfLaTeX (biber automático).
% Guía completa: apa7/README.md
\usepackage[style=apa, backend=biber, sortcites=true, sorting=nyt]{biblatex}
\DeclareLanguageMapping{spanish}{spanish-apa}
\addbibresource{referencias.bib}   % copie aquí su .bib (modelo: apa7/referencias_apa.bib)

% ============================================
% CONFIGURACIÓN DE MÁRGENES (ESTILO UNTELS)
% ============================================
\geometry{
    left=4.0cm,    % Espacio para encuadernación
    right=2.5cm,
    top=2.5cm,
    bottom=2.5cm
}

% ============================================
% CONFIGURACIÓN DE ESPACIADO
% ============================================
\onehalfspacing  % Espaciado 1.5 líneas

% ============================================
% CONFIGURACIÓN DE ENCABEZADOS
% ============================================
\pagestyle{fancy}
\fancyhf{}
\fancyhead[L]{\leftmark}
\fancyhead[R]{\thepage}
\renewcommand{\headrulewidth}{0.4pt}
\renewcommand{\chaptermark}[1]{\markboth{\chaptername\ \thechapter.\ #1}{}}

% ============================================
% CONFIGURACIÓN DE TÍTULOS DE CAPÍTULOS
% ============================================
\titleformat{\chapter}
    {\normalfont\LARGE\bfseries\centering}
    {\chaptertitlename\ \thechapter}{1em}{}
\titlespacing*{\chapter}{0pt}{-20pt}{20pt}

% ============================================
% METADATOS DE LA TESIS
% ============================================
\title{Título de la Tesis}
\author{Apellidos y Nombres del Autor}
\date{\today}

% ============================================
% INICIO DEL DOCUMENTO
% ============================================
\begin{document}

% ============================================
% PORTADA
% ============================================
\begin{titlepage}
    \centering
    \vspace*{1cm}
    
    \includegraphics[width=0.3\textwidth]{logo_untels.png}
    \vspace{1cm}
    
    {\large \textbf{UNIVERSIDAD NACIONAL TECNOLÓGICA DE LIMA SUR}}
    \vspace{0.3cm}
    
    {\large FACULTAD DE INGENIERÍA ELECTRÓNICA}
    \vspace{2cm}
    
    {\Large \textbf{TESIS}}
    \vspace{0.5cm}
    
    {\LARGE \textbf{Título completo de la tesis de investigación}}
    \vspace{0.3cm}
    
    {\large Presentada para optar el título profesional de:}
    \vspace{0.3cm}
    
    {\large \textbf{INGENIERO ELECTRÓNICO}}
    \vspace{2cm}
    
    {\large \textbf{Autor:}}
    \vspace{0.2cm}
    
    {\large Apellidos y Nombres del Autor}
    \vspace{0.5cm}
    
    {\large \textbf{Asesor:}}
    \vspace{0.2cm}
    
    {\large Nombre del Asesor, Grado Académico}
    \vspace{1cm}
    
    {\large Lima, Perú}
    \vspace{0.2cm}
    
    {\large \today}
\end{titlepage}

% ============================================
% DEDICATORIA
% ============================================
\newpage
\thispagestyle{empty}
\vspace*{8cm}
\begin{center}
    \textit{A mis padres, por su apoyo incondicional} \\
    \textit{A mis profesores, por guiar mi formación} \\
    \textit{A Dios, por darme fortaleza}
\end{center}

% ============================================
% AGRADECIMIENTOS
% ============================================
\newpage
\chapter*{AGRADECIMIENTOS}
\addcontentsline{toc}{chapter}{AGRADECIMIENTOS}

Agradezco a la Universidad Nacional Tecnológica de Lima Sur por brindarme la formación académica y profesional. A mi asesor, [Nombre], por su guía y paciencia durante el desarrollo de esta investigación. A mis compañeros y amigos, por su apoyo y motivación constante.

% ============================================
% DECLARACIÓN DE USO DE HERRAMIENTAS DE IA
% ============================================
% Obligatoria en esta plantilla. Si NO usó IA, deje solo el primer párrafo
% alternativo y borre el resto. Más detalle: etica/README.md
\newpage
\chapter*{DECLARACIÓN DE USO DE HERRAMIENTAS DE INTELIGENCIA ARTIFICIAL}
\addcontentsline{toc}{chapter}{DECLARACIÓN DE USO DE HERRAMIENTAS DE INTELIGENCIA ARTIFICIAL}

% --- Opción A: SÍ se usó IA (adapte a lo que realmente hizo) ---
Durante la elaboración de esta tesis, el autor utilizó [herramienta, modelo/versión, proveedor] entre [fecha de inicio] y [fecha de fin] con la finalidad de [finalidad concreta: depurar código \LaTeX, mejorar la redacción de la sección X, etc.]. Tras su uso, el autor revisó, verificó y editó el contenido generado y asume la plena responsabilidad por el trabajo presentado.

\textbf{Datos y resultados.} Todos los datos experimentales fueron obtenidos y procesados por el autor; ningún dato ni resultado fue generado por inteligencia artificial. [Si hubo datos simulados: los datos de la Tabla [N.°] son simulados y así se rotulan en el texto.]

\textbf{Referencias.} Todas las fuentes citadas fueron consultadas directamente por el autor; ninguna fue incorporada a partir de una lista generada por IA sin verificación.

\textbf{Trazabilidad.} El registro de las consultas realizadas (fecha, modelo, \textit{prompt}, verificación y cambios) se presenta en el Apéndice~\ref{apx:registro-ia} y los \textit{prompts} y respuestas completos se conservan en [ubicación] a disposición del asesor y del jurado.

% --- Opción B: NO se usó IA (borre la Opción A y use solo esto) ---
% Durante la elaboración de esta tesis el autor no utilizó herramientas de inteligencia artificial generativa.

\vspace{1cm}
\begin{flushright}
    [Lugar], [fecha] \\[1.5cm]
    \rule{6cm}{0.4pt} \\
    {[Nombre y apellidos del autor]} \\
    DNI {[N.°]}
\end{flushright}

% ============================================
% RESUMEN
% ============================================
\newpage
\chapter*{RESUMEN}
\addcontentsline{toc}{chapter}{RESUMEN}

Este trabajo de tesis presenta el desarrollo e implementación de [descripción breve del proyecto]. El objetivo principal es [objetivo general]. La metodología empleada incluye [metodología]. Los resultados obtenidos demuestran [resultados principales]. Se concluye que [conclusiones].

\textbf{Palabras clave:} palabra1, palabra2, palabra3, palabra4, palabra5.

\vspace{1cm}

\chapter*{ABSTRACT}
\addcontentsline{toc}{chapter}{ABSTRACT}

This thesis presents the development and implementation of [brief description]. The main objective is [general objective]. The methodology includes [methodology]. The results demonstrate [main results]. It is concluded that [conclusions].

\textbf{Keywords:} keyword1, keyword2, keyword3, keyword4, keyword5.

% ============================================
% ÍNDICES
% ============================================
\newpage
\tableofcontents
\newpage
\listoffigures
\newpage
\listoftables

% ============================================
% CAPÍTULO 1: INTRODUCCIÓN
% ============================================
\chapter{INTRODUCCIÓN}

\section{Planteamiento del Problema}
Describir el problema que motiva la investigación.

\subsection{Descripción de la realidad problemática}
\textbf{Contenido:} Explicar la situación actual, identificar carencias o dificultades.

\subsection{Definición del problema}
\textbf{Contenido:} Formular el problema de investigación de manera clara y precisa.

\subsection{Justificación e importancia}
\textbf{Contenido:} Explicar por qué es relevante realizar esta investigación.

\section{Objetivos de la Investigación}

\subsection{Objetivo General}
Desarrollar un [sistema/método/prototipo] para [aplicación específica].

\subsection{Objetivos Específicos}
\begin{enumerate}
    \item Objetivo específico 1
    \item Objetivo específico 2
    \item Objetivo específico 3
    \item Objetivo específico 4
\end{enumerate}

\section{Hipótesis}
\textbf{Hipótesis general:} [Formular la hipótesis principal].

\textbf{Hipótesis específicas:}
\begin{itemize}
    \item H1: [Hipótesis específica 1]
    \item H2: [Hipótesis específica 2]
    \item H3: [Hipótesis específica 3]
\end{itemize}

\section{Alcance y Limitaciones}
\subsection{Alcance}
Describir los límites y cobertura de la investigación.

\subsection{Limitaciones}
Identificar las restricciones encontradas durante el desarrollo.

% ============================================
% CAPÍTULO 2: MARCO TEÓRICO
% ============================================
\chapter{MARCO TEÓRICO}

\section{Antecedentes de la Investigación}
\subsection{Investigaciones previas nacionales}
\subsection{Investigaciones previas internacionales}

\section{Bases Teóricas}
\subsection{Principios fundamentales}
La presente investigación se fundamenta en los principios de la electrónica analógica y digital \parencite{boylestad2022}.
% Citas APA 7: \parencite{clave} -> (Autor, año) | \textcite{clave} -> Autor (año)
% | \parencite[p.~25]{clave} -> (Autor, año, p. 25). No use \cite (sin paréntesis).

\subsection{Leyes y teoremas}
La Ley de Ohm establece la relación entre voltaje, corriente y resistencia:

\begin{equation}
    V = I \cdot R
    \label{eq:ohm}
\end{equation}

\subsection{Modelos matemáticos}
El comportamiento de un circuito RC se describe mediante:

\begin{equation}
    V_c(t) = V_{in} \left(1 - e^{-t/\tau}\right)
    \label{eq:rc_carga}
\end{equation}

donde $\tau = RC$ es la constante de tiempo del circuito.

\section{Marco Conceptual}
\subsection{Definición de términos básicos}
\begin{itemize}
    \item \textbf{Voltaje:} Diferencia de potencial eléctrico entre dos puntos.
    \item \textbf{Corriente:} Flujo de carga eléctrica por unidad de tiempo.
    \item \textbf{Resistencia:} Oposición al flujo de corriente eléctrica.
\end{itemize}

\section{Estado del Arte}
\subsection{Tecnologías existentes}
\subsection{Tendencias actuales}

% ============================================
% CAPÍTULO 3: METODOLOGÍA
% ============================================
\chapter{METODOLOGÍA}

\section{Tipo y Nivel de Investigación}
\subsection{Tipo de investigación}
Aplicada / Básica / Tecnológica.

\subsection{Nivel de investigación}
Descriptiva / Correlacional / Explicativa.

\section{Diseño de la Investigación}
\subsection{Diseño experimental / no experimental}
\subsection{Esquema del diseño}

\section{Población y Muestra}

\section{Técnicas e Instrumentos de Recolección de Datos}
\subsection{Técnicas empleadas}
\subsection{Instrumentos utilizados}

\begin{table}[H]
    \centering
    \caption{Instrumentos de medición utilizados}
    \begin{tabular}{@{}lll@{}}
        \toprule
        \textbf{Instrumento} & \textbf{Marca} & \textbf{Precisión} \\
        \midrule
        Osciloscopio & Tektronix & $\pm$0.1\% \\
        Multímetro & Fluke & $\pm$0.5\% \\
        Fuente DC & GwInstek & $\pm$0.01\% \\
        \bottomrule
    \end{tabular}
    \label{tab:instrumentos}
\end{table}

\section{Uso de Herramientas de Inteligencia Artificial}
% Complete SOLO si la IA influyó en el método (análisis de datos, código, simulación).
% Para el uso editorial o de formato basta la Declaración inicial.
Las herramientas de IA empleadas, su finalidad y la forma de verificación de sus resultados se describen a continuación: [herramienta y versión]; [tarea: p.\,ej. script de ajuste de curvas]; [verificación: p.\,ej. se comprobaron dos puntos manualmente y se contrastó con la simulación]. El código generado se incluye en el Apéndice [X] y fue ejecutado con los datos medidos por el autor.

\section{Procedimiento Experimental}
\subsection{Etapa 1: Diseño y simulación}
\subsection{Etapa 2: Implementación}
\subsection{Etapa 3: Pruebas y mediciones}

% ============================================
% CAPÍTULO 4: DESARROLLO E IMPLEMENTACIÓN
% ============================================
\chapter{DESARROLLO E IMPLEMENTACIÓN}

\section{Diseño del Sistema}
\subsection{Diagrama de bloques}
\begin{figure}[H]
    \centering
    \framebox[0.8\textwidth]{\rule{0pt}{4cm} \textbf{Diagrama de bloques}}
    \caption{Diagrama de bloques del sistema propuesto}
    \label{fig:bloques}
\end{figure}

\subsection{Especificaciones técnicas}
\begin{table}[H]
    \centering
    \caption{Especificaciones del sistema}
    \begin{tabular}{@{}ll@{}}
        \toprule
        \textbf{Parámetro} & \textbf{Valor} \\
        \midrule
        Voltaje de alimentación & 5 V DC \\
        Consumo de corriente & 150 mA \\
        Frecuencia de operación & 16 MHz \\
        Resolución ADC & 10 bits \\
        \bottomrule
    \end{tabular}
    \label{tab:especificaciones}
\end{table}

\section{Implementación del Hardware}
\subsection{Selección de componentes}
\subsection{Diagrama esquemático}
\subsection{Diseño de PCB}

\section{Implementación del Software}
\subsection{Algoritmo principal}
\begin{lstlisting}[language=C, caption=Código del sistema]
#include <stdio.h>
#include <stdlib.h>

int main() {
    // Configuración inicial
    printf("Sistema iniciado\n");
    
    // Bucle principal
    while(1) {
        // Lectura de sensores
        // Procesamiento de datos
        // Control de actuadores
    }
    return 0;
}
\end{lstlisting}

\subsection{Interfaz de usuario}
\subsection{Protocolos de comunicación}

% ============================================
% CAPÍTULO 5: RESULTADOS Y DISCUSIÓN
% ============================================
\chapter{RESULTADOS Y DISCUSIÓN}

\section{Presentación de Resultados}

\subsection{Resultados experimentales}
\begin{table}[H]
    \centering
    \caption{Resultados de las mediciones realizadas}
    \begin{tabular}{@{}cccc@{}}
        \toprule
        \textbf{Prueba} & \textbf{Valor teórico} & \textbf{Valor medido} & \textbf{Error (\%)} \\
        \midrule
        1 & 5.00 V & 4.98 V & 0.4 \\
        2 & 10.00 V & 9.95 V & 0.5 \\
        3 & 15.00 V & 14.92 V & 0.53 \\
        \bottomrule
    \end{tabular}
    \label{tab:resultados}
\end{table}

\subsection{Análisis estadístico}
\subsection{Gráficas comparativas}

\section{Discusión de Resultados}
\subsection{Interpretación de resultados}
\subsection{Comparación con trabajos previos}
\subsection{Aportes de la investigación}

\section{Validación de Hipótesis}
\subsection{Verificación de hipótesis general}
\subsection{Verificación de hipótesis específicas}

% ============================================
% CAPÍTULO 6: CONCLUSIONES Y RECOMENDACIONES
% ============================================
\chapter{CONCLUSIONES Y RECOMENDACIONES}

\section{Conclusiones}
\begin{enumerate}
    \item Conclusión derivada del objetivo específico 1
    \item Conclusión derivada del objetivo específico 2
    \item Conclusión derivada del objetivo específico 3
    \item Conclusión derivada del objetivo específico 4
    \item Conclusión general derivada del objetivo general
\end{enumerate}

\section{Recomendaciones}
\begin{enumerate}
    \item Recomendación para futuras investigaciones
    \item Recomendación para aplicaciones prácticas
    \item Recomendación para mejora del sistema
\end{enumerate}

\section{Aportes de la Investigación}
\subsection{Aportes teóricos}
\subsection{Aportes prácticos}
\subsection{Aportes metodológicos}

% ============================================
% REFERENCIAS BIBLIOGRÁFICAS
% ============================================
% Los títulos de libros y revistas en cursiva, la sangría francesa y el orden
% alfabético los aplica biblatex-apa automáticamente.
\printbibliography[title={REFERENCIAS BIBLIOGRÁFICAS}, heading=bibintoc]

% ============================================
% APÉNDICES
% ============================================
\begin{appendices}
\chapter{REGISTRO DE USO DE HERRAMIENTAS DE IA}
\label{apx:registro-ia}
% Una fila por uso relevante. Guarde también los prompts completos (Git, carpeta, etc.).
\begin{longtable}{>{\raggedright}p{2.6cm}>{\raggedright}p{3cm}>{\raggedright}p{3.4cm}>{\raggedright\arraybackslash}p{5.2cm}}
    \caption{Registro de consultas a herramientas de IA}\label{tab:registro-ia}\\
    \toprule
    \textbf{Fecha y modelo} & \textbf{Parte de la tesis} & \textbf{Prompt (resumen)} & \textbf{Verificación y cambios realizados} \\
    \midrule
    \endfirsthead
    \toprule
    \textbf{Fecha y modelo} & \textbf{Parte de la tesis} & \textbf{Prompt (resumen)} & \textbf{Verificación y cambios realizados} \\
    \midrule
    \endhead
    \bottomrule
    \endfoot
    {[AAAA-MM-DD, modelo]} & {[Capítulo/sección]} & {[Resumen del prompt]} & {[Qué comprobó y qué cambió o rechazó]} \\
    {[\dots]} & {[\dots]} & {[\dots]} & {[\dots]} \\
\end{longtable}

\chapter{MANUAL DE USUARIO}
Contenido del manual de usuario del sistema desarrollado.

\chapter{DATOS TÉCNICOS}
Fichas técnicas de los componentes utilizados.

\chapter{CÓDIGO FUENTE}
Listado completo del código fuente implementado.

\end{appendices}

\end{document}
